\subsection{结构画像：为什么需要多种初解}
在建立具体求解器之前，可用操作数、弱连通分量数、最大分量规模与最长拓扑深度描述图的\emph{结构}，但不能仅凭这些静态量断言最终周期。表\ref{tab:six-structure}列出本工程早期用于开发与确认的六张图的结构快照；它们不是由最终性能反推而后选出的“六张有利图”，也不是一百图的全体统计。可切计算操作数与总操作数不同，因为原始 COPY 等节点不由外层分配子图。最大分量占比 $\chi$ 按式\eqref{eq:q2-dominant-ratio}用可切操作数计算。第 $48$ 图近乎单一深长依赖分量，第 $18$、$33$ 图却有大量小分量；对这两类图使用相同的“每块固定算子数”会在一类图上切碎关键链，在另一类图上错失天然独立块。因此分量完整性、主导分量有界切分和拓扑窗口必须作为不同初解来源并存。
\begin{table}[!htbp]\centering\small
\caption{六张结构诊断图的静态画像；不是全百图的分布统计}\label{tab:six-structure}
\begin{tabular}{lrrrrr}
\toprule
图号 & 总操作 & 可切计算操作 & 弱连通分量 & 最大分量 & 最大拓扑深度\\
\midrule
12 & 1304 & 972 & 50 & 98 & 18\\
48 & 1900 & 1728 & 1 & 1728 & 368\\
50 & 4781 & 4555 & 23 & 4011 & 185\\
18 & 3285 & 2727 & 126 & 59 & 27\\
33 & 4069 & 3465 & 95 & 64 & 33\\
67 & 8996 & 8804 & 71 & 124 & 123\\
\bottomrule
\end{tabular}
\end{table}

具体计算最大分量占比：第 $48$ 图 $\chi=1728/1728=1$；第 $50$ 图 $\chi=4011/4555\approx0.8806$。完整分量直接分核，至少有一个核心需要承担主导分量的大部分工作，难以充分利用其余核心。相反，第 $18$ 图 $\chi=59/2727\approx0.0216$，第 $33$ 图 $\chi=64/3465\approx0.0185$，分量整体分核已经提供很多可调的独立工作块；对它们过早细切单个小分量，很可能增加 Task 或跨核边界而无明显并行收益。第 $67$ 图操作多，却有 $71$ 个分量、最大分量仅 $124$ 个可切操作，提示“图很大”不等于“需要把某一个分量切成极小块”。这说明主机运行耗时与 NPU 最优划分尺度是不同问题：第 $67$ 图因为节点多、官方核内展开复杂而求解贵，仍可先尝试以完整分量为单位构造高质量初解。

这些结构量也并非独立决定瓶颈。最大拓扑深度是按操作层最长链的节点数，忽略每个操作的周期、Pipe 类别与张量字节；两条深度相同的链可以有完全不同的计算时长和 DDR 压力。弱连通分量不计外部输入共享，因此不同分量即使没有中间依赖，也可能共同消费同一输入张量，分给不同 Task 后重复读取。最大分量占比则只用算子数量，没有权衡某些算子的高周期或大张量。正因如此，表\ref{tab:six-structure}只用于初解门控和结构解释，不能替代官方评价或作为全局最优的预言变量。

\subsection{静态结构与五核实测的交叉核对}
为检验结构判断有没有和实际结果完全脱节，表\ref{tab:six-outcomes}把同六张图在五核下的两种场景实测并列。第 $48$ 图在 A 中只取得约 $1.1245$ 倍，在 B 中提高到约 $2.1863$ 倍；第 $50$ 图从约 $1.5062$ 倍到约 $2.6254$ 倍。它们的主导分量和长链为“完整分量单核”带来明显限制，场景 B 允许同核驻留并重构候选后表现更好。不过这个成对结果同时改变了 Task 规则、同步、正式候选集，不能把差值单独归因于某一个因素。第 $18$、$33$ 图无主导分量，A 的五核表现已接近五倍；B 中绝对周期反而略高。第 $67$ 图 A/B 五核周期也几乎持平。这些反例说明后续场景的机制更丰富，不意味着每张图都严格比前一问快。
\begin{table}[!htbp]\centering\small
\caption{结构诊断图的五核实测对照；两场景使用各自已选方案}\label{tab:six-outcomes}
\begin{tabular}{lrrrr}
\toprule
图号 & A 周期 & A 逐图加速比 & B 周期 & B 逐图加速比\\
\midrule
12 & 13406 & 4.4076 & 12519 & 4.7199\\
48 & 197613 & 1.1245 & 101641 & 2.1863\\
50 & 99382 & 1.5062 & 57017 & 2.6254\\
18 & 203806 & 4.9281 & 205917 & 4.8775\\
33 & 353592 & 5.0176 & 361176 & 4.9123\\
67 & 13455316 & 4.7442 & 13461153 & 4.7422\\
\bottomrule
\end{tabular}
\end{table}

对第 $48$ 图，仅看“最大分量占比为一”还不足以决定把链切在何处；若切点位于共享张量的高扇出处，会产生大量边界 COPY。可沿关键输出的祖先链选择候选切点，估计切开后双计算 Pipe 的并行余量，再查看跨切点张量集合和释放点。对第 $18$、$33$ 图，完整分量数量已远超过五核，最重要的问题转为把分量按两 Pipe 负载放到核心、以及控制共享输入重复读取；继续细化每个很小的分量大概率只是增加候选准备成本。对第 $67$ 图，五核周期很大却已有超过四倍的相对加速，若以赛题逐图均值为目标，它贡献与任一小图相同；若以总周期为目标，它的绝对改善又非常重要。结构分型因此同时服务于\emph{哪类初解有价值}和\emph{主机搜索时间投向哪里}两项不同决策。

\subsection{从三问差异推出统一而不混淆的求解路线}
共同变量 $p=(\pi,m,\sigma)$ 允许三个问题共享合法性框架和结构候选，但不允许共用一个不变的通信权重矩阵。场景 A 的同核跨 Task 张量也必须经 DDR，故主要边界是子图；场景 B 的同核多个子图在同一 Task 中可驻留，故主要边界转成核心；场景 C 又使相同 B 方案中的某些读入由 L2 而非 DDR 服务，并引入全局 FIFO 历史。若按“通信成本从 A 到 C 逐渐减小”简单描述，会漏掉 L2 竞争与非单调反例。严谨的递进关系是三个评价算子
\begin{equation}\label{eq:three-evaluators}
\mathfrak E_A\circ\mathfrak K\circ\mathfrak G_A,
\quad
\mathfrak E_B\circ\mathfrak K\circ\mathfrak G_B,
\quad
\mathfrak E_C\circ\mathfrak K\circ\mathfrak G_B,
\end{equation}
其中 $\mathfrak E_C$ 比 $\mathfrak E_B$ 多共享 L2 队列和第二服务池，而 $\mathfrak G_A$ 与 $\mathfrak G_B$ 采用不同 Task 与 COPY 规则。这个表达把可复用的\emph{图结构候选}与必须重评的\emph{物理成本}分开，解释了为何问题二以问题一方案作种子、却全部重放正式评价；问题三以问题二方案作保底、再做固定方案消融与少量缓存事件候选。

算法最终回答的是每个配置的具体方案，而非仅给出理论表达式。对一张图，结构初解多样化是为了避免落入单一分块尺度；关键工作量和空隙插入是为了让已存在的计算资源不被低优先级任务占住；有限官方评价是为了让 DDR 竞争、核内溢出与 FIFO 路由这些非线性因素进入最终接受判断。若后续证明某个解析代理完全复刻官方仿真，它在计算成本上就应按正式评价记账；在本论文中代理始终被限定为有误差的排序工具。统一框架的优势不是宣称一个权重公式适用于三问，而是在每问按相应事件语义重新评估的同时保持候选结构、合法性和证据口径的一致。
